\subsection{Wrapping Level}
\label{sec:wrapping}

The wrapping level $w$ classifies every UI element into one of four structural categories based on its boundary semantics:

\begin{center}
\begin{tabular}{cl}
\toprule
$w$ & Boundary class \\
\midrule
0 & Inline / no boundary \\
1 & Single-line bounded control \\
2 & Multi-line bounded block \\
3 & Structural section / large overlay \\
\bottomrule
\end{tabular}
\end{center}

\noindent Representative examples for each level:

\begin{itemize}
  \item $w = 0$: plain text, icon, inline label. No outer boundary.
  \item $w = 1$: button, text input, select, tooltip. A single-line container with visible or implicit edges.
  \item $w = 2$: textarea, blockquote, card. A multi-line region with a defined boundary.
  \item $w = 3$: dialog, drawer, fieldset. A structural overlay or section that frames nested content.
\end{itemize}

The wrapping level is an intrinsic property of the element's role, not a runtime parameter. It is assigned once per patch type and does not change with density or font size.

This classification is exhaustive: every element in the UI library maps to exactly one wrapping level. Section~\ref{sec:dimension} confirms this by tabulating all production patches.
